\documentclass[UTF8,11pt]{ctexart}
\usepackage{amsmath,amssymb,geometry,hyperref}
\geometry{a4paper,margin=25mm}
\title{经典自旋的谱矩与低阶 Krylov 递推}
\date{}
\begin{document}
\maketitle
\noindent\textbf{形式化范围：}谱矩积分的闭式表达已写成 Lean 定理；
完整经典自旋模型和任意阶 Lanczos 系数尚未形式化。编译状态以仓库 CI 为准。

对固定长度经典自旋 $S$，考虑 $H=\chi S_z^2$ 和观测量 $A=S_x$。
令 $z=S_z/S$、$c=2\chi S$。由经典运动方程并对球面均匀测度
$d\mu=dz\,d\phi/(4\pi)$ 取平均，归一化关联函数可写为
\begin{equation}
 C(t)=\frac34\int_{-1}^1(1-z^2)\cos(czt)\,dz.
 \label{eq:corr}
\end{equation}
其对应的偶数谱矩定义为
\begin{equation}
 \mu_{2n}=\frac34c^{2n}\int_{-1}^1 (1-z^2)z^{2n}\,dz.
\end{equation}
利用定积分线性性和幂积分公式，
\[
 \begin{aligned}
 \int_{-1}^1(1-z^2)z^{2n}\,dz
 &=\int_{-1}^1z^{2n}\,dz-
 \int_{-1}^1z^{2n+2}\,dz\\
 &=\frac2{2n+1}-\frac2{2n+3}
 =\frac4{(2n+1)(2n+3)}.
 \end{aligned}
\]
因此
\begin{equation}
 \boxed{\mu_{2n}=\frac{3c^{2n}}{(2n+1)(2n+3)}}.
 \label{eq:moments}
\end{equation}
式 \eqref{eq:moments} 的\textbf{积分数学命题}在
\texttt{ClassicalSpinKrylov.evenMoment\_formula} 中形式化。
特别地 $\mu_0=1$，$\mu_2=c^2/5$，$\mu_4=3c^4/35$。

利用零对角 Lanczos 链的矩关系（该关系目前尚未形式化）
$\mu_2=b_1^2$、$\mu_4=b_1^4+b_1^2b_2^2$，在 $c\ne0$ 时
\[
 b_1^2=\frac{c^2}{5},\qquad
 b_2^2=\frac{\mu_4-\mu_2^2}{\mu_2}=\frac{8c^2}{35}.
\]
两项低阶结果的\textbf{代数一致性}分别由
\texttt{fourth\_moment\_identity} 和
\texttt{second\_lanczos\_identity} 检查。
完整的 Lanczos 通项还需要另外证明，不能从本文已验证结果直接推出。
\end{document}
